\documentclass[11pt, a4paper]{article}

% bfw-style.tex — gemeinsame Style-Definitionen für alle Handouts/Aufgabenhefte
% Voraussetzung: Kompilation mit xelatex (wegen fontspec)

\usepackage[ngerman]{babel}
\usepackage{geometry}
\geometry{a4paper, top=2.2cm, bottom=2.2cm, left=2.3cm, right=2.3cm}

\usepackage{fontspec}
% Fallback-Kette: Source Sans Pro -> Calibri -> Segoe UI -> Latin Modern Sans
% (Latin Modern ist immer mit MiKTeX vorhanden)
\IfFontExistsTF{Source Sans Pro}{%
  \setmainfont{Source Sans Pro}[Ligatures=TeX]%
}{\IfFontExistsTF{Calibri}{%
  \setmainfont{Calibri}[Ligatures=TeX]%
}{\IfFontExistsTF{Segoe UI}{%
  \setmainfont{Segoe UI}[Ligatures=TeX]%
}{%
  \setmainfont{Latin Modern Sans}[Ligatures=TeX]%
}}}
\IfFontExistsTF{Source Code Pro}{%
  \setmonofont{Source Code Pro}[Scale=0.92]%
}{\IfFontExistsTF{Consolas}{%
  \setmonofont{Consolas}[Scale=0.92]%
}{%
  \setmonofont{Latin Modern Mono}[Scale=0.92]%
}}

\usepackage{xcolor}
% Farbpalette: hell, freundlich, modern
\definecolor{bfwPrimary}{HTML}{0EA5A4}    % Petrol/Türkis - Primär
\definecolor{bfwPrimaryDark}{HTML}{0F766E} % dunkler Petrol für Headings
\definecolor{bfwAccent}{HTML}{F59E0B}     % Amber - Akzent
\definecolor{bfwSuccess}{HTML}{10B981}    % Grün - Erfolg / Lösung
\definecolor{bfwWarn}{HTML}{F97316}       % Orange - Achtung
\definecolor{bfwInk}{HTML}{0F172A}        % fast Schwarz
\definecolor{bfwInkSoft}{HTML}{475569}    % gedämpftes Grau
\definecolor{bfwBgSoft}{HTML}{F8FAFC}     % off-white
\definecolor{bfwBgTeal}{HTML}{ECFEFF}     % helles Türkis
\definecolor{bfwBgAmber}{HTML}{FEF3C7}    % helles Gelb
\definecolor{bfwBgRose}{HTML}{FFE4E6}     % helles Rosé für Eselsbrücken

\usepackage[colorlinks=true, linkcolor=bfwPrimaryDark, urlcolor=bfwPrimaryDark]{hyperref}
\usepackage{enumitem}
\setlist[itemize]{leftmargin=*, itemsep=2pt, topsep=2pt}
\setlist[enumerate]{leftmargin=*, itemsep=2pt, topsep=2pt}

\usepackage[most]{tcolorbox}
\usepackage{booktabs}
\usepackage{tikz}
\usetikzlibrary{decorations.pathmorphing, positioning, shapes.geometric, arrows.meta}

% Merkkästchen — wichtige Definition / Kernpunkt
\newtcolorbox{merksatz}[1][]{%
  enhanced, breakable,
  colback=bfwBgTeal,
  colframe=bfwPrimary,
  boxrule=0.6pt,
  arc=4pt,
  left=10pt, right=10pt, top=8pt, bottom=8pt,
  fonttitle=\bfseries\color{bfwPrimaryDark},
  title={\faLightbulb\ Merksatz},
  #1
}

% Eselsbrücke — Mnemoniks
\newtcolorbox{eselsbruecke}[1][]{%
  enhanced, breakable,
  colback=bfwBgRose,
  colframe=bfwAccent,
  boxrule=0.6pt,
  arc=4pt,
  left=10pt, right=10pt, top=8pt, bottom=8pt,
  fonttitle=\bfseries\color{bfwWarn},
  title={Eselsbrücke},
  #1
}

% Beispiel-Box
\newtcolorbox{beispiel}[1][]{%
  enhanced, breakable,
  colback=bfwBgSoft,
  colframe=bfwInkSoft,
  boxrule=0.4pt,
  arc=3pt,
  left=10pt, right=10pt, top=8pt, bottom=8pt,
  fonttitle=\bfseries\color{bfwInk},
  title={Beispiel},
  #1
}

% Lösungs-Box (für Aufgabenhefte)
\newtcolorbox{loesung}[1][]{%
  enhanced, breakable,
  colback=bfwBgAmber!40,
  colframe=bfwSuccess,
  boxrule=0.6pt,
  arc=3pt,
  left=10pt, right=10pt, top=8pt, bottom=8pt,
  fonttitle=\bfseries\color{bfwSuccess!70!black},
  title={Lösung},
  #1
}

% Glossar-Eintrag
\newcommand{\glossareintrag}[2]{%
  \par\noindent\textbf{\color{bfwPrimaryDark}#1} \enspace #2\par\smallskip
}

% Section-Stil — etwas farbiger als Standard
\usepackage{titlesec}
\titleformat{\section}
  {\Large\bfseries\color{bfwPrimaryDark}}
  {\thesection}{0.6em}{}
\titleformat{\subsection}
  {\large\bfseries\color{bfwPrimary}}
  {\thesubsection}{0.5em}{}
\titleformat{\subsubsection}
  {\normalsize\bfseries\color{bfwInk}}
  {\thesubsubsection}{0.4em}{}

% TikZ-Stil "skizzenhaft" — etwas weicher als Rough.js, aber stabil
\tikzset{
  roughbox/.style = {
    draw=bfwPrimaryDark, thick, fill=bfwBgTeal,
    rounded corners=4pt, inner sep=6pt
  },
  roughaccent/.style = {
    draw=bfwAccent, thick, fill=bfwBgAmber,
    rounded corners=4pt, inner sep=6pt
  },
  roughline/.style = {
    draw=bfwInkSoft, thick, -{Stealth[length=2.5mm]}
  }
}

% Bullet-Symbole (bewusst kein FontAwesome — vermeidet Font-Installations-Probleme)
\newcommand{\faLightbulb}{$\bullet$}
\newcommand{\faBookmark}{$\bullet$}

% Header/Footer
\usepackage{fancyhdr}
\pagestyle{fancy}
\fancyhf{}
\renewcommand{\headrulewidth}{0pt}
\fancyfoot[L]{\small\color{bfwInkSoft}\thepage}
\fancyfoot[R]{\small\color{bfwInkSoft} BFW M-064 Beschaffung im Büromanagement}


\title{\vspace*{-1.5cm}\Huge\bfseries\color{bfwPrimaryDark}Aufgabenheft Tag~1\\[0.4em]
  \large\color{bfwInkSoft}\textnormal{Büroraumplanung \& Büroformen --- Übungen im IHK-Klausurformat}}
\author{\textsf{Büroprozesse gestalten und Arbeitsvorgänge organisieren \textperiodcentered{} M-067}\\
\textsf{\small\copyright\ Can Siebert\,\textperiodcentered\,2026}}
\date{Selbstlernmaterial \textperiodcentered{} 29.06.2026}

\begin{document}
\maketitle
\thispagestyle{empty}

\begin{merksatz}[title={\bfseries So nutzen Sie dieses Aufgabenheft}]
Bearbeiten Sie alle Aufgaben zunächst \emph{ohne} in die Lösungen zu schauen. Schreiben
Sie Ihre Antworten in vollständigen Sätzen --- die IHK-Prüfung verlangt formulierte
Begründungen, nicht bloße Stichworte. Beachten Sie die \emph{Operatoren} (nennen,
erläutern, beurteilen, begründen) --- sie geben den geforderten Antwort-Tiefegrad vor.
Erst danach öffnen Sie den Lösungsteil ab Seite~\pageref{loesungen} und vergleichen.
\end{merksatz}

\tableofcontents
\vspace{1em}
\hrule

\section{Aufgaben}

\subsection*{Aufgabe~1 --- Bedeutung der Büroraumplanung (10 Punkte)}

\noindent\textbf{\color{bfwAccent}YouTube-Vorbereitung:}\enspace \href{https://www.youtube.com/results?search_query=Bueroraumplanung+Grundlagen}{Büroraumplanung --- Grundlagen} (\url{https://www.youtube.com/results?search_query=Bueroraumplanung+Grundlagen})

\medskip
Sie sind in der Verwaltung der \emph{Lindner \& Partner GmbH} tätig. Die Geschäfts\-führung
plant einen Umzug und bittet Sie, die Bedeutung einer sorgfältigen Büroraumplanung zu
erläutern.

\begin{enumerate}[label=(\alph*)]
  \item \textbf{Erläutern} Sie, warum die Raumgestaltung eines Büros für ein Unternehmen wichtig ist. \hfill (3~P.)
  \item \textbf{Nennen} Sie drei Faktoren, die durch gute Büroraumplanung positiv beeinflusst werden. \hfill (3~P.)
  \item \textbf{Beschreiben} Sie kurz, wofür die DGUV Information 215-441 dient. \hfill (2~P.)
  \item \textbf{Beurteilen} Sie die Aussage: \emph{„Ein schöneres Büro ist reine Geldverschwendung.”} \hfill (2~P.)
\end{enumerate}

\subsection*{Aufgabe~2 --- Die sechs Teilflächen (12 Punkte)}

\noindent\textbf{\color{bfwAccent}YouTube-Vorbereitung:}\enspace \href{https://www.youtube.com/results?search_query=Flaechenarten+Bueroarbeitsplatz}{Flächenarten am Büroarbeitsplatz} (\url{https://www.youtube.com/results?search_query=Flaechenarten+Bueroarbeitsplatz})

\medskip
Die Fläche eines Büros gliedert sich in sechs Teilflächen.

\begin{enumerate}[label=(\alph*)]
  \item \textbf{Nennen} Sie alle sechs Teilflächen eines Büros. \hfill (3~P.)
  \item \textbf{Ordnen} Sie jeder Teilfläche ein konkretes Praxisbeispiel \textbf{zu}. \hfill (3~P.)
  \item \textbf{Erläutern} Sie den Unterschied zwischen \emph{Funktionsfläche} und \emph{Bewegungsfläche}. \hfill (3~P.)
  \item \textbf{Beurteilen} Sie, warum man den Flächenbedarf nicht einfach pauschal schätzen sollte. \hfill (3~P.)
\end{enumerate}

\subsection*{Aufgabe~3 --- Flächenbedarf berechnen (12 Punkte)}

\noindent\textbf{\color{bfwAccent}YouTube-Vorbereitung:}\enspace \href{https://www.youtube.com/results?search_query=Quadratmeter+pro+Arbeitsplatz+Buero}{Quadratmeter pro Arbeitsplatz} (\url{https://www.youtube.com/results?search_query=Quadratmeter+pro+Arbeitsplatz+Buero})

\medskip
Ein Unternehmen plant eine neue Büroetage. Für einen Bildschirm\-arbeitsplatz werden
8--10~m\textsuperscript{2} veranschlagt, für einen Arbeitsplatz im Großraumbüro
ca.\ 12--15~m\textsuperscript{2}.

\begin{enumerate}[label=(\alph*)]
  \item \textbf{Berechnen} Sie, wie viele Bildschirm\-arbeitsplätze zu je 10~m\textsuperscript{2} in einen Raum von 80~m\textsuperscript{2} passen. \hfill (3~P.)
  \item \textbf{Berechnen} Sie, wie viele Großraum\-arbeitsplätze zu je 12~m\textsuperscript{2} in dieselben 80~m\textsuperscript{2} passen, und \textbf{vergleichen} Sie. \hfill (3~P.)
  \item \textbf{Erläutern} Sie, warum ein Arbeitsplatz im Großraumbüro mehr Fläche benötigt als ein einzelner Bildschirm\-arbeitsplatz. \hfill (3~P.)
  \item \textbf{Erklären} Sie, wie sich der Quadratmeter\-bedarf pro Arbeitsplatz grundsätzlich ergibt. \hfill (3~P.)
\end{enumerate}

\subsection*{Aufgabe~4 --- Büroformen vergleichen (15 Punkte)}

\noindent\textbf{\color{bfwAccent}YouTube-Vorbereitung:}\enspace \href{https://www.youtube.com/results?search_query=Bueroformen+Zellenbuero+Grossraumbuero}{Büroformen im Vergleich} (\url{https://www.youtube.com/results?search_query=Bueroformen+Zellenbuero+Grossraumbuero})

\medskip
Die \emph{Lindner \& Partner GmbH} (Steuerberatung, 18 Beschäftigte) zieht um. Es gibt
konzentrierte Einzelarbeit (Mandantenakten), Teamarbeit (Projekte) und vertrauliche
Mandantengespräche. Die Geschäftsführung diskutiert zwischen Zellen-, Gruppen- und
Kombibüro.

\begin{enumerate}[label=(\alph*)]
  \item \textbf{Nennen} Sie je zwei Vor- und Nachteile von Zellen-, Gruppen- und Großraumbüro. \hfill (6~P.)
  \item \textbf{Begründen} Sie, welche Büroform Sie der Steuerberatung empfehlen. \hfill (5~P.)
  \item \textbf{Nennen} Sie zwei Maßnahmen, mit denen sich der Geräuschpegel in einem Gruppenbüro senken lässt. \hfill (2~P.)
  \item \textbf{Beurteilen} Sie: Warum gilt das Großraumbüro trotz Kostenvorteils oft als problematisch? \hfill (2~P.)
\end{enumerate}

\subsection*{Aufgabe~5 --- Moderne Bürokonzepte (12 Punkte)}

\noindent\textbf{\color{bfwAccent}YouTube-Vorbereitung:}\enspace \href{https://www.youtube.com/results?search_query=Desksharing+non+territoriales+Buero}{Desksharing \& non-territoriales Büro} (\url{https://www.youtube.com/results?search_query=Desksharing+non+territoriales+Buero})

\medskip
Ein Unternehmen mit 100 Beschäftigten stellt fest, dass durch Außendienst, Projektarbeit
und Homeoffice täglich nur etwa 60 Personen gleichzeitig im Büro sind.

\begin{enumerate}[label=(\alph*)]
  \item \textbf{Erläutern} Sie das Prinzip des \emph{non-territorialen Büros} (Desksharing). \hfill (3~P.)
  \item \textbf{Nennen} Sie zwei Vorteile und einen Nachteil dieses Konzepts. \hfill (3~P.)
  \item \textbf{Erläutern} Sie den Grundgedanken des \emph{reversiblen Büros}. \hfill (3~P.)
  \item \textbf{Beurteilen} Sie, wie sich zunehmendes Homeoffice auf den Flächenbedarf des Unternehmens auswirkt. \hfill (3~P.)
\end{enumerate}

\subsection*{Aufgabe~6 --- Büroform situativ zuordnen (10 Punkte)}

\noindent\textbf{\color{bfwAccent}YouTube-Vorbereitung:}\enspace \href{https://www.youtube.com/results?search_query=Kombibuero+Gruppenbuero+Unterschied}{Kombibüro vs.\ Gruppenbüro} (\url{https://www.youtube.com/results?search_query=Kombibuero+Gruppenbuero+Unterschied})

\medskip
\begin{enumerate}[label=(\alph*)]
  \item \textbf{Charakterisieren} Sie kurz die vier klassischen Büroformen (Zellen-, Gruppen-, Großraum-, Kombibüro). \hfill (4~P.)
  \item \textbf{Ordnen} Sie für jeden Fall die passende Büroform zu \textbf{und begründen} Sie: \hfill (4~P.)
    \begin{enumerate}[label=\arabic*., leftmargin=*]
      \item Ein Call-Center mit 120 Telefon-Arbeitsplätzen soll möglichst kostengünstig untergebracht werden.
      \item Eine Geschäftsführerin braucht einen Raum für vertrauliche Gespräche und konzentriertes Arbeiten.
      \item Ein Projektteam von acht Personen muss sich ständig abstimmen.
      \item Eine Abteilung braucht störungsfreie Einzelarbeit \emph{und} gelegentliche Teamarbeit.
    \end{enumerate}
  \item \textbf{Beurteilen} Sie das größte \emph{Problem} des Großraumbüros anhand eines Beispiels. \hfill (2~P.)
\end{enumerate}

\vspace{1em}
\noindent\textbf{Punkte gesamt: 71}

\newpage

\section{Lösungen}\label{loesungen}

\subsection*{Lösung Aufgabe~1}
\begin{loesung}
\textbf{(a)} Die Raumgestaltung eines Büros wirkt sich entscheidend auf die Produktivität
und die Arbeitszufriedenheit der Nutzer aus. Da Büroflächen außerdem erhebliche Kosten
(Miete, Energie, Reinigung) verursachen, sollten die räumlichen Ressourcen sorgfältig
geplant werden --- gute Planung ist also zugleich Wirtschaftlichkeit und Personalfürsorge.

\textbf{(b)} Drei positiv beeinflusste Faktoren (Beispiele): Produktivität,
Arbeitszufriedenheit, Gesundheit/Sicherheit, Konzentrationsfähigkeit, Kommunikation.

\textbf{(c)} Die DGUV Information 215-441 „Büroraumplanung” dient als Leitfaden. Sie
konkretisiert die sicherheits\-technischen, arbeitsmedizinischen, ergonomischen und
arbeits\-psychologischen Anforderungen an die Planung von Büroflächen.

\textbf{(d)} Die Aussage ist falsch. Ein gut geplantes Büro ist keine Verschwendung,
sondern eine Investition: Zufriedene, gesunde Beschäftigte sind produktiver, seltener
krank und bleiben dem Unternehmen länger erhalten. Schlechte Gestaltung verursacht
Lärm, Krankheit und höhere Fehlzeiten --- und damit am Ende höhere Kosten.
\end{loesung}

\subsection*{Lösung Aufgabe~2}
\begin{loesung}
\textbf{(a)} Die sechs Teilflächen: \textbf{Stellfläche, Bewegungsfläche, Funktionsfläche,
Benutzerfläche am Arbeitsplatz, Verkehrswegefläche, Besucher-/Besprechungsfläche.}

\textbf{(b)} Beispiele:
\begin{itemize}[itemsep=0pt]
  \item \emph{Stellfläche}: Bodenfläche unter Schreibtisch, Schrank, Regal, Pflanze.
  \item \emph{Bewegungsfläche}: freier Bereich am Stuhl zum Aufstehen/Drehen.
  \item \emph{Funktionsfläche}: Bereich vor dem Schrank, der beim Öffnen der Tür frei sein muss.
  \item \emph{Benutzerfläche}: Fläche am persönlichen Arbeitsplatz für die Tätigkeit.
  \item \emph{Verkehrswegefläche}: Gang zwischen den Arbeitsplätzen.
  \item \emph{Besucher-/Besprechungsfläche}: Besucherstuhl, kleiner Besprechungstisch.
\end{itemize}

\textbf{(c)} Die \emph{Funktionsfläche} wird beim Öffnen von Schranktüren oder Schubladen
\emph{überdeckt} --- sie hängt also an den Möbeln. Die \emph{Bewegungsfläche} ist
dagegen freie, unverstellte Fläche am Arbeitsplatz, die dem Menschen wechselnde
Arbeitshaltungen und Ausgleichs\-bewegungen ermöglicht.

\textbf{(d)} Der Flächenbedarf entsteht aus der \emph{Addition} der sechs Teilflächen.
Eine pauschale Schätzung würde einzelne Flächen (z.\,B.\ Funktions- oder Verkehrswege\-fläche)
vergessen --- die Folge wären zu enge, unsichere oder ineffiziente Arbeitsplätze.
\end{loesung}

\subsection*{Lösung Aufgabe~3}
\begin{loesung}
\textbf{(a)} $80~\text{m}^2 \div 10~\text{m}^2 = \textbf{8 Bildschirmarbeitsplätze}$.

\textbf{(b)} $80~\text{m}^2 \div 12~\text{m}^2 = 6{,}67$, also \textbf{6 Großraum\-arbeitsplätze}
(volle Plätze). Im Vergleich passen bei Großraum-Maßstab \emph{weniger} Plätze auf
dieselbe Fläche, weil je Platz mehr Fläche eingeplant wird.

\textbf{(c)} Im Großraumbüro sind mehr \emph{Verkehrswege} notwendig: Je mehr Menschen
sich einen Raum teilen, desto mehr Fläche entfällt anteilig auf Wege zwischen den
Arbeitsplätzen. Deshalb 12--15~m\textsuperscript{2} statt 8--10~m\textsuperscript{2}.

\textbf{(d)} Der Quadratmeter\-bedarf pro Arbeitsplatz ergibt sich aus der \emph{Addition
der jeweils ermittelten Werte der sechs Teilflächen}.
\end{loesung}

\subsection*{Lösung Aufgabe~4}
\begin{loesung}
\textbf{(a)} Vor-/Nachteile:
\begin{itemize}[itemsep=0pt]
  \item \emph{Zellenbüro}: $+$ konzentriertes/vertrauliches Arbeiten, individuelle Gestaltung; $-$ Isolationsgefühl, hoher Flächenbedarf.
  \item \emph{Gruppenbüro}: $+$ intensive Kommunikation, schneller Informationsfluss, leichte Vertretung; $-$ hoher Geräuschpegel, weniger Privatsphäre.
  \item \emph{Großraumbüro}: $+$ kostengünstig, flexibel, informeller Wissensaustausch; $-$ mangelnde Privatsphäre, Kontrollgefühl, höherer Krankenstand.
\end{itemize}

\textbf{(b)} Empfehlung: \textbf{Kombibüro}. Es vereint die Vorteile von Zellen- und
Großraumbüro: Die Einzelbüros an der Fensterfront ermöglichen störungsfreie Einzelarbeit
und vertrauliche Mandantengespräche, die zentrale Multifunktionszone bietet Raum für
Teamarbeit und Projektbesprechungen. Das passt genau zur Mischung aus Konzentration,
Vertraulichkeit und Kooperation in der Steuerberatung.

\textbf{(c)} Maßnahmen gegen Geräuschpegel: Telefon- und Arbeitszellen als
Rückzugsmöglichkeit, schallabsorbierende Decken, schalldämpfende Teppiche.

\textbf{(d)} Das Großraumbüro ist zwar kostengünstig (effiziente Flächennutzung, weniger
Geräte), doch die Ersparnis wird laut Studien häufig durch Unzufriedenheit, höheren
Krankenstand und geringere Produktivität wieder aufgehoben --- der Kostenvorteil ist
also oft nur scheinbar.
\end{loesung}

\subsection*{Lösung Aufgabe~5}
\begin{loesung}
\textbf{(a)} Beim non-territorialen Büro ist die Anzahl der Arbeitsplätze \emph{kleiner}
als die Zahl der Beschäftigten. Nicht jeder hat einen persönlichen Schreibtisch; bei
Arbeitsbeginn sucht man sich den nächsten freien Platz (oft per Buchungssoftware).
Persönliche Unterlagen werden in einem fahrbaren Rollcontainer aufbewahrt, fachliche
Unterlagen digital oder zentral abgelegt.

\textbf{(b)} Vorteile: bessere Auslastung der Fläche, erhebliche Kosteneinsparung
(Büroflächen, Möbel, Technik, Reinigung, Energie). Nachteil: keine feste Anlaufstelle
für Kunden/Kollegen, und der Arbeitsplatz muss jeweils neu an die Körpergröße angepasst
werden.

\textbf{(c)} Das reversible Büro kann mit \emph{geringem Aufwand} an organisatorische
Erfordernisse angepasst werden. Statt fester Wände werden flexible Wandsysteme genutzt;
so kann jede Etage anforderungsgerecht eine andere Büroform (Zellen-, Gruppen-,
Kombibüro) aufnehmen.

\textbf{(d)} Bei 100 Beschäftigten, aber nur ca.\ 60 gleichzeitig Anwesenden werden
nicht 100 Arbeitsplätze benötigt. Homeoffice senkt also den gleichzeitigen
Flächenbedarf; Desksharing wird wirtschaftlich attraktiv, und das Unternehmen kann
Mietfläche und Folgekosten einsparen --- vorausgesetzt, klare Regeln und verlässliche
Technik sind vorhanden.
\end{loesung}

\subsection*{Lösung Aufgabe~6}
\begin{loesung}
\textbf{(a)} Charakterisierung:
\begin{itemize}[itemsep=0pt]
  \item \emph{Zellenbüro}: 1--3 Personen, Ruhe, Vertraulichkeit, individuelle Gestaltung.
  \item \emph{Gruppenbüro}: 3--25 Plätze, Kommunikation und Teamarbeit, aber laut.
  \item \emph{Großraumbüro}: 25 bis mehrere hundert Plätze, kostengünstig, wenig Privatsphäre.
  \item \emph{Kombibüro}: Einzelbüros plus Multifunktionszone, vereint Einzel- und Teamarbeit.
\end{itemize}

\textbf{(b)} Zuordnung mit Begründung:
\begin{enumerate}[label=\arabic*., itemsep=0pt]
  \item Call-Center: \textbf{Großraumbüro} --- kostengünstig durch effiziente Flächennutzung; typischer Einsatzbereich.
  \item Geschäftsführerin: \textbf{Zellenbüro} --- vertrauliche Gespräche, konzentriertes Arbeiten, Statussymbol.
  \item Projektteam: \textbf{Gruppenbüro} --- fördert ständige Abstimmung, Kommunikation und kurzen Informationsweg.
  \item Einzel- und Teamarbeit: \textbf{Kombibüro} --- Einzelbüros für Konzentration, Multifunktionszone für Teamarbeit.
\end{enumerate}

\textbf{(c)} Größtes Problem des Großraumbüros: \emph{mangelnde Privatsphäre und
Geräuschpegel}. Beispiel: In einem Großraum mit 60 Telefon-Arbeitsplätzen stören
Gespräche und Telefonate die Konzentration ständig; das Gefühl ständiger Kontrolle und
fehlende Rückzugsmöglichkeiten führen zu Stress, höherem Krankenstand und sinkender
Produktivität.
\end{loesung}

\end{document}
